\documentclass[12pt]{article}
\usepackage{amsmath, amssymb, enumitem, fancyhdr, graphicx, libertine, nth,
  pgfplots, tikz, xcolor}
\pgfplotsset{compat=1.11}
\usetikzlibrary{decorations.pathreplacing}
\usetikzlibrary{patterns}
\usepackage[margin=1in]{geometry}
\usepackage[libertine]{newtxmath}
\pagestyle{fancy}
\definecolor{solution-color}{HTML}{000080}
\chead{\emph{EC201 -- Problem Set 5 -- Solutions}}

% Define new topic command
\newcommand{\topic}[1]{\medskip\begin{center}\large\textsc{#1}\normalsize
  \end{center}}

% Define question counter:
\newcounter{question}[section]
\newcommand{\question}[1]{\refstepcounter{question}
  \subsubsection*{Question~\thequestion~-- #1}}

% Subquestions are enumerate with letters
\newenvironment{subquestions}{
  \begin{enumerate}[label=(\roman*)]}{\end{enumerate}}

\begin{document}
\topic{Price Discrimination}

\question{Third-Degree Price Discrimination}
There is one movie theater in a town. The town has two groups of people:
students and non-students. The non-students have a demand curve:
\[
  D_1\left( p_1\right) = 10-p_1
\]
and the students have a demand curve:
\[
  D_2\left( p_2 \right) = 10 - 2p_2
\]
The movie theater has a marginal cost of \$2 per customer (cleaning up their
spilled soda and popcorn). The fixed cost is 10.
\begin{subquestions}
  \item Suppose first the monopolist movie theater cannot price discriminate.
    What price should they charge and what quantity will it sell? What
    profit does it make?

    \textcolor{solution-color}{
      The full demand curve is:
      \[
        D\left( p\right)=D_1\left( p \right)+D_2\left( p \right)=10-p_1 + 10 -
        2p = 20-3p
      \]
      The inverse demand curve is then:
      \[
        p=\frac{20}{3}-\frac{q}{3}
      \]
      The revenue function is then:
      \[
        R\left( q \right)=p\left( q \right)q=\frac{20}{3}q -
        \frac{1}{3}q^2
      \]
      The marginal revenue is then:
      \[
        MR\left( q \right)=\frac{20}{3}-\frac{2}{3}q
      \]
      Setting $MR\left( q \right)=MC\left( q \right)$ gives:
      \[
        \frac{20}{3}-\frac{2}{3}q=2 \;\;\;\;\Longrightarrow \;\;\;\;
        q = 7
      \]
      The price is
      $\frac{20}{3}-\frac{7}{3}=\frac{13}{3}=4\frac{1}{3}$.
      The profit is then:
      \[
        \pi=\frac{13}{3}\times 7  - 2\times 7 - 10 = 6\frac{1}{3}
      \]
    }
  \item Suppose now the monopolist is able to price discriminate. What price
    should it charge students and what price should it charge non-students?
    How many tickets will they sell?
    \textcolor{solution-color}{
      The individual inverse demand curves are:
      \[
        p_1\left( q_1 \right)=10-q_1
      \]
      and:
      \[
        p_2\left( q_2 \right)=5-\frac{1}{2}q_2
      \]
      The revenue and marginal revenue functions for each group are:
      \[
        R_1\left( q_1 \right)=10q_1-q_1^2 \;\;\;\;\;\;\;\;\;\;\;\;\;\;
        MR_1\left( q_1 \right)=10-2q_1
      \]
      and:
      \[
        R_2\left( q_2 \right)=5q_2-\frac{1}{2}q^2_2 \;\;\;\;\;\;\;\;\;\;\;\;\;\;
        MR_2\left( q_2 \right)=5 - q_2
      \]
      The monopolist's problem is then to set the marginal revenue from each
      group equal to marginal cost:
      \[
        10-2q_1=2 \;\;\;\;\; \Longrightarrow\;\;\;\;\;  q_1 = 4
      \]
      \[
        5-q_2=2 \;\;\;\;\; \Longrightarrow\;\;\;\;\;  q_2 = 3
      \]
      The prices are:
      \[
        p_1=10-4=6 \;\;\;\;\;\;\;\;\;\;\;\; p_2 =
        5-\frac{3}{2}=3\frac{1}{2}
      \]
      The profit is then:
      \[
        \pi= p_1q_1 + p_2q_2 - \left( q_1+q_2 \right)c - F=6\times 4 + \left(
        3\frac{1}{2} \right)\times 3 - \left( 3+4 \right)\times 2 - 10=10.5
      \]
      The firm makes a higher profit by by discriminiating.
    }
\end{subquestions}
\question{Monopoly Price Discrimination}
A monopolist airline faces two types of consumers who value flight quality
differently. There is a high willingness-to-pay group (business people) and a
low willingness-to-pay group (vacationers). The airline faces zero marginal
cost (putting one more person on a plane doens't any extra).

``Quality" is measured on a 30-point scale, where 30 is extremely
luxurious and 0 is extremely uncomfortable.

The two solid downward-sloping lines below measure the willingness-to-pay for a
marginal increase in quality for both groups.  For any quality level, $q$, the
willingness-to-pay is the area \emph{under the line} between between 0 and $q$. Each triangle blow has an area of 50 and each square has an area of 100.
If $q=10$, vacationers are willing to pay at most $150$ for air travel,
while business people would pay $250$.  If the quality level is 0,
the area is always 0 so no one would pay for airfare at quality 0.
\begin{center}
  \begin{tikzpicture}[scale=0.8]
    \draw[thick,<->] (0,9) node[above] {$WTP$}--(0,0)--(9,0) node[below right]
    {Quality, $q$};
    \draw[dashed,-] (0,5)--(2.5,5);
    \draw[dashed,-] (5,0)--(5,2.5);
    \draw[dashed,-] (2.5,0)--(2.5,5);
    \draw[dashed,-] (0,2.5)--(5,2.5);
    \draw (0, 7.5) node[left] {30};
    \draw (0, 5) node[left] {20};
    \draw (0, 2.5) node[left] {10};
    \draw (2.5, 0) node[below] {10};
    \draw (5, 0) node[below] {20};
    \draw (7.5, 0) node[below] {30};
    \draw (0, 0) node[below left] {0};
    \draw (0,7.5)--(7.5,0) node[above right] {\textsl{Business people}};
    \draw (0,5)--(5,0) node[above right] {\textsl{Vacationers}};
  \end{tikzpicture}
\end{center}
\textcolor{solution-color}{We can summarize the willingness to pay for each type for each quality level as follows:
  \begin{center}
    \begin{tabular}{l|l|l} \hline \hline
      & Vacationers & Business people \\ \hline
      Low quality ($q=10$)  &  150 & 250 \\
      Medium quality ($q=20$) & 200 & 400 \\
      High quality ($q=30$) & 200 & 450 \\ \hline \hline
    \end{tabular}
  \end{center}
}
\begin{subquestions}
  \item If the only people in the world were business people, what should the airline charge for air travel to maximize profits?

    \textcolor{solution-color}{
      You should set the quality level to 30 and charge exactly what the market will value the flight. This is the area of the entire triangle which is 450.
    }
  \item If the only people in the world were vacationers, what should the airline charge for air travel to maximize profits?

    \textcolor{solution-color}{
      You should set the quality level to 20 and charge exactly what the market will value the flight. This is the area of the smaller triangle which is 200.
    }
\end{subquestions}
For the rest of the question, assume that 50\% of air travelers are business
people and 50\% are vacationers.
\begin{subquestions}
  \setcounter{enumi}{2}
  \item Suppose vacationers were able to obtain a ``vacationer's card'' which allowed them to get discounts for air travel if discounts were available (like a student card). The airline now has an observable way to tell travelers apart. What should the airline charge each group to maximize profits? What is the average profit made per person?\footnote{If you get revenue $x$ from vactioners and revenue $y$ from business people, the average profit per person is $\frac{x+y}{2}$. If you only sell to business people, the average profit per person is $\frac{y}{2}$.}

      \textcolor{solution-color}{
        The airline should just charge each group for the different quality levels in (i) and (ii). The vacationers wouldn't want to buy the high quality flight (they would get negative surplus by doing so) so they wouldn't pretend to be business people. Business people would want to go on the lower-quality flight, but since they don't have ``vacationers'' cards they can't get tickets for it.  Therefore the airline offers a high quality flight with $q_H=30$ for 450 and a medium-quality flight with $q_M=20$ for 200. The average profit made per person is $0.5\times 450 + 0.5\times 200= 325$.
      }
  \item Suppose now that a vacationer's card doesn't exist so the airline can't charge different people different prices.  If the airline were to choose only one price-quality plan, what would it do to maximize profits? To do this you should check two possibilities:
    \begin{itemize}
      \item You set a price-quality plan such that both groups would be
        willing to buy.
      \item You set a price-quality plan such that only business people will
        be willing to buy (but you lose out on half the market).
    \end{itemize}
    Find the average profit made per person from both of these possibilities.

      \textcolor{solution-color}{
        \begin{itemize}
          \item If the airline wanted to sell to both groups it should offer $q=20$ for 200. Vacationers would buy it and get zero surplus. The business people would buy it and get surplus equal to 200. The average profit per person is 200 (since everyone buys it).
        \item If the airline focused only on the business people, the airline should offer $q=30$ like in (i) for 450. The average profit per person is then $0.5\times 450=225$ (since you lose the vacationers).
        \item Therefore if the airline was only going to set a single price-quality pair, it should sell only to business people.
      \end{itemize}
      }
  \item Suppose again that the vacationer's card doesn't exist but the airline now wants to set two price-quality plans in order to make the vacationers and business people \emph{self-select} into each plan.
    \begin{itemize}
      \item What price-quality plans for both groups would make the groups self-select and maximize profits?
      \item What is the average profit per customer?
    \end{itemize}

      \textcolor{solution-color}{
        Here we need to check the different combinations of quality levels and choose the one that maximizes profits.
      }

      \textcolor{solution-color}{
        We should also check the other combinations of qualities to confirm that this indeed is the profit-maximizing strategy:
        \begin{itemize}
          \item $q_L=10$ and $q_M=30$. If the airline tried to sell $q=10$ to vacationers it would sell it for 150 (their willingness to pay). Business people would get 100 surplus from this so they need to give business people 100 surplus on the high quality flight. Their willingness to pay for the high quality flight is 450 so they should charge them 350 to ensure they are not tempted to buy the low quality flight. The average profit made per person is then $0.5\times 150+0.5 \times 350=250$. 
          \item $q_M=20$ and $q_H=30$. If the airline tried to sell $q=20$ to vacationers it would sell it for 200. Business people would get 200 surplus from this so they should charge 450-200=250 for the $q_H=30$ flight. On average this earns $0.5\times 200 + 0.5 \times 250=225$ which is less than 250 from before.
          \item $q_L=10$ and $q_M=20$. If the airline tried to sell $q=10$ to vacationers it would sell it for 150. Business people would get 100 surplus from this so they should charge 400-100=300 for the $q_H=20$ flight. On average this earns $0.5\times 150 + 0.5 \times 300=225$ which is less than 250 from before.
        \end{itemize}
        Therefore the best plan for the airline is to offer a low quality flight $q_L=10$ for 150 and a high quality flight $q_H=30$ for 450. The airline makes a higher profit with two different options than by offering only one, even though it can't tell the two groups of consumers apart.
      }
\end{subquestions}
% \question{Bundling}
% You are a monopolist software company sell three different software products.
% Three types of people buy your software: Marketers, Accountants and Graphic
% Designers. There is an equal amount of each type of person in the population.
% The willingness to pay of each group for each type software is below:
% \begin{center}
%   \begin{tabular}{|l|c|c|c|}\hline
%                            & Marketers   & Accountants & Graphic Designers \\
%     \hline
%     Presentation Software  & 150         & 100         & 0                 \\
%     Spreadsheet Software   & 100         & 150         & 0                 \\
%     Image-Editing Software & 40          & 10          & 150               \\
%     \hline
%   \end{tabular}
% \end{center}
% You are not able to price discriminate between the different groups.
% \begin{subquestions}
%   \item If you were selling each piece of software individually, what price
%     should you charge for each?

%     \textcolor{solution-color}{
%       If you sold the presentation software for \$0, everyone would buy it. If
%       you sold it for \$100, $\frac{2}{3}$ of the market would buy it (so the
%       average revenue would be $66\frac{2}{3}$). If you sold it for \$150, only
%       $\frac{1}{3}$ of the market would buy it (so the average revenue would be
%       \$50). Therefore you should charge \$100 for the presentation software.
%     }

%     \textcolor{solution-color}{
%       Using this approach for the other two pieces of software, you will decide to
%       charge \$100 for the spreadsheet software and \$150 for the image-editing
%       software.  In total, the average revenue from all 3 softwares is
%       $66\frac{2}{3}+66\frac{2}{3}+50=183.34$.
%     }
%   \item If you were able to bundle two or three pieces of software together,
%     what pieces of software would you bundle together?

%     \textcolor{solution-color}{
%       If you bundled the presentation and spreadsheet software together you
%       could sell it for \$250 and make an average revenue of \$166.67 on those
%       two pieces of software (previously you sold each of them for \$100 each,
%       making an average revenue of $\$66\frac{2}{3}$ each). Selling the
%       image-editing software then for \$150 would finally leave you with an
%       average profit of \$166.67+\$50=\$216.67.
%     }

%     \textcolor{solution-color}{
%       You wouldn't want to sell all three individually as we saw in part (i)
%       that the average revenue is only \$183.34.
%     }

%     \textcolor{solution-color}{
%         You wouldn't want to bundle the 3
%         pieces of software together. The graphic designers value all 3 together
%         the least at $0+0+150=150$. If you sold the bundle for \$150,
%         everyone would buy it so the average revenue would be \$150.
%         If you sold all three for \$260, the marketers and accountants would
%         buy it but the average revenue would be only \$173.33. At any
%         higher price for all three, the average revenue would be even less.
%       }

%     \textcolor{solution-color}{
%       Therefore the best bundling strategy is to bundle the presentation and
%       speadsheet software together and sell the image-editing software on its
%       own and get \$216.67 on average.
%       }

%     \textcolor{solution-color}{
%       There is another profitable bundling choice you could do. If you were to
%       sell all three as a bundle \emph{and} charge separately for the image
%       editing software, You could charge \$260 for all three piece of software,
%       and sell to $\frac{2}{3}$ of the market (making on average \$173.33).
%       Then you could also sell the image-editing software for \$150 on its own
%       (making on average \$50). This together would give \$223.33 on average.
%     }
% \end{subquestions}


\question{Two-Part Tariffs}
There is a single theme park in a large state. It can put people on a roller coaster at zero marginal cost (but they do have a fixed cost). The demand for roller coaster rides for each person is given by $D\left( p \right)=200-\frac{1}{2}p$, where $p$ is the price of each roller coaster ride. What price should the theme park charge for entry into the theme park and how much should it charge for each person to ride the roller coaster?

\textcolor{solution-color}{
  We saw in class the optimal thing to do in a two-part tariff is to charge at
  marginal cost (here zero) and then have the entrance fee equal to the
  consumer surplus that the consumers would get if you only charged at marginal
  cost.
  The inverse demand curve is $p\left( q \right)=400-2q$. The area of the
  triangle under this demand curve is $\frac{1}{2}\times 400 \times 200 = 40,000$
  Therefore the theme park should charge an entry fee of \$40,000 and charge
  nothing for the rides. This steals the entire consumer surplus from consumers
  and there is no deadweight loss.
}

\topic{Game Theory}
\question{Simulataneous Games}
Find all Nash equilibria (pure and mixed) of the following games:

  \textcolor{solution-color}{
    In the following, for mixed-strategy equilibria, $p$ is the probability
  player 1 goes up ($U$) and $q$ is the probability that player 2 goes left
  ($L$).
  }
\begin{subquestions}
  \item
  \setlength{\arraycolsep}{0pt}
  \[
  \begin{array}{cccccc}
  &  & \quad L \quad & & \quad R \quad  &  \\
  \cline{3-6} U \quad & \vline{} & \quad 2,4  \quad & \vline{} & \quad 6,3 \quad & \vline{}  \\
  \cline{3-6} D \quad & \vline{} & \quad 3,1 \quad & \vline{} & \quad 5,3 \quad & \vline{} \\
  \cline{3-6}
  \end{array}
  \]
  \textcolor{solution-color}{
    There are no pure strategy equilibria. There is a mixed-strategy
    equilibrium with $p=\frac{2}{3}$ and $q=\frac{1}{2}$. To see how we get these probabilities, player 1 plays $U$ with probability $p$ and player 2 plays $L$ with probability $q$. For player 1 to be indifferent between the two actions we need:
    \begin{align*}
  2q+6\left( 1-q \right)&=3q+5\left( 1-q \right) \\
  2q + 6 - 6q &= 3q+5-5q\\
  1&=2q\\
  q=\frac{1}{2}
    \end{align*}
    For player 2 to be indifferent between the two actions we need:
    \begin{align*}
      4p+\left( 1-p \right)&=3p+3\left( 1-p \right)\\
      4p+1-p&=3\\
      3p&=2\\
      p=\frac{2}{3}
    \end{align*}
  }

  \item
  \setlength{\arraycolsep}{0pt}
  \[
  \begin{array}{cccccc}
  &  & \quad L \quad & & \quad R \quad  &  \\
  \cline{3-6} U \quad & \vline{} & \quad 2,1  \quad & \vline{} & \quad 0,0 \quad & \vline{}  \\
  \cline{3-6} D \quad & \vline{} & \quad 0,0 \quad & \vline{} & \quad 1,2 \quad & \vline{} \\
  \cline{3-6}
  \end{array}
  \]

  \textcolor{solution-color}{
    $(U,L)$ and $(D,R)$ are pure strategy equilibria. There is also a mixed
    strategy equilibrium with $p=\frac{2}{3}$ and $q=\frac{1}{3}$.
  }

  \item
  \setlength{\arraycolsep}{0pt}
  \[
  \begin{array}{cccccc}
  &  & \quad L \quad & & \quad R \quad  &  \\
  \cline{3-6} U \quad & \vline{} & \quad -2,-2  \quad & \vline{} & \quad 4,0 \quad & \vline{}  \\
  \cline{3-6} D \quad & \vline{} & \quad 0,4 \quad & \vline{} & \quad 2,2 \quad & \vline{} \\
  \cline{3-6}
  \end{array}
  \]

  \textcolor{solution-color}{
    $(D,L)$ and $(U,R)$ are pure strategy equilibria. There is also a mixed
    strategy equilibrium with $p=q=\frac{1}{2}$.
}
\end{subquestions}

\question{Sequential Games}
Consider the following extensive form game:
\ \\
\begin{center}
\begin{tikzpicture}[scale=1.5,font=\footnotesize]
\tikzstyle{solid node}=[circle,draw,inner sep=1.5,fill=black]
\tikzstyle{hollow node}=[circle,draw,inner sep=1.5]
\tikzstyle{level 1}=[level distance=15mm,sibling distance=3.5cm]
\tikzstyle{level 2}=[level distance=15mm,sibling distance=1.5cm]
\tikzstyle{level 3}=[level distance=15mm,sibling distance=1cm]
\node(0)[hollow node,label=above:{$P1$}]{}
child{node[solid node,label=above left:{$P2$}]{}
child{node[label=below:{$(5,3)$}]{} edge from parent node[left]{$A$}}
child{node[label=below:{$(0,2)$}]{} edge from parent node[left]{$B$}}
edge from parent node[left,xshift=-5]{$L$}
}
child{node[label=below:{$(4,2)$}]{} edge from parent node [above]{$R$}
};
\end{tikzpicture}
\end{center}
\begin{subquestions}
  \item Write the normal form (payoff matrix) of the game.

    \textcolor{solution-color}{
      The normal form is:
      \setlength{\arraycolsep}{0pt}
      \[
      \begin{array}{cccccc} &  & \quad A \quad & & \quad B \quad  &  \\
        \cline{3-6} L \quad & \vline{} & \quad 5,3  \quad & \vline{} & \quad
        0,2 \quad & \vline{}  \\ \cline{3-6} R \quad & \vline{} & \quad 4,2
        \quad & \vline{} & \quad 4,2 \quad & \vline{} \\ \cline{3-6}
      \end{array}
      \]
    }
  \item Find all the Nash equilibria of the game.

    \textcolor{solution-color}{
      There are two pure strategy NE: $(R,B)$ and $(L,A)$. There is also a MSNE where P1 does $R$ always and P2 plays $A$ with probability $\frac{4}{5}$ and $B$ with probability $\frac{1}{5}$, i.e.~$p=0$ and $q=\frac{4}{4}$.
    }
  \item Which of the Nash equilibria that you found are subgame perfect?

    \textcolor{solution-color}{
        $(L,A)$ is the only SPNE.
    }
\end{subquestions}
\newpage
\topic{Oligopoly}
\question{Different Market Structures}
The demand curve for a good is given by $D\left( p
\right)=5-\frac{1}{2}p$. The cost function for a firm is given by $c\left( q
\right)=2q$.
\begin{subquestions}
  \item The market was served by perfectly competitive firms.
      What is the market price and total output produced by the industry?  What
      profit does each firm make?

      \textcolor{solution-color}{
        In perfect competition, price equals marginal cost. The marginal cost
        here is $c^\prime\left( q \right)=2$. With $p=2$ we can use the demand
        function to get $q=5-\frac{1}{2}\times 2 = 4$. Each firm makes zero
        profit since price equals marginal cost and there is no fixed cost.
      }
  \item The market is served by a monopolist. What price and quantity does the
    monopolist produce? What profit does it make?

      \textcolor{solution-color}{
        The inverse demand curve is $p\left( q \right)=10-2q$. The revenue
        function is then $R\left( q \right)=10q - 2q^2$. The marginal revenue
        function is then $MR\left( q \right)=R^\prime\left( q \right)=10-4q$.
        In monopoly, to maximize profits you set $MR\left( q \right)=MC\left(
        q \right)$ to find the optimal quantity. So:
        \[
          10-4q=2 \;\;\;\;\Longleftrightarrow \;\;\;\; q=2
        \]
        Note that this is half of the perfectly competitive quantity.
        The price then is $p=10-2\times 2 = 6$. The monopolist's profits is
        then $\pi=\left(p-AC  \right)q = \left( 6-2 \right)2=8$.
      }
  \item The market is served by two firms who compete simultaneously on price.
    What price does each firm charge? What quantity is produced by both firms?
    What profit does each firm make?

    \textcolor{solution-color}{
      This is Bertrand competition. If two firms compete on price then they
      will engage in a price war and end up charging $p=MC$, like in perfect
      competition. So $p=2$, the quantity produced in total will be $q=4$ and
      both firms will make zero profits.
    }
  \item  The market is served by two firms who simulatenously choose their
    quantites. The price is then set to clear the market. What quantity is
    produced by both firms? What is the price? What is the profit for each
    firm?

    \textcolor{solution-color}{
      This is Cournot competition.
      In general terms, firm 1's problem is:
      \[
        \underset{q_1}{\max}\; p\left( q_1+q_2 \right)q_1 - c\left( q_1
        \right)
      \]
      The revenue term is:
      \[
        p\left( q_1+q_2 \right)q_1=\left[ 10-2\left( q_1+q_2
      \right) \right]q_1= 10q_1 - 2q_1^2 -2q_1q_2
      \]
      Therefore firm 1's problem is:
      \[
        \underset{q_1}{\max}\;  10q_1 - 2q_1^2 -2q_1q_2-2q_1
      \]
      The first-order condition is then:
      \[
        10 - 4q_1 - 2q_2 - 2 =0
      \]
      So $q_1\left( q_2 \right) = 2 - \frac{1}{2}q_2$. This is firm 1's best
      response function.
      For firm 2, we can do the same steps and arrive at $q_2\left( q_1
      \right) =
      2-\frac{1}{2}q_1$.
      In equilibrium, both equations will be satisfied:
      \begin{align*}
        q_1&=2-\frac{1}{2}q_2 \\
        q_2&= 2-\frac{1}{2}q_1
      \end{align*}
      Inserting the equation for $q_2$ into the equation for $q_1$:
      \begin{align*}
        q_1 &= 2-\frac{1}{2}\left( 2 - \frac{1}{2}q_1 \right) \\
        q_1 &= 2 - 1 + \frac{1}{4}q_1 \\
        \frac{3}{4}q_1 &= 1 \\
        q_1 &= \frac{4}{3}
      \end{align*}
      Using this in the best response function for firm 2 we get:
      $q_2=2-\frac{1}{2}\times \frac{4}{3}
      =2-\frac{4}{6}=\frac{4}{3}$ also. Each firm produces one third of
      the perfectly competitive quantity. The industry as a whole then
      produces $q=2\times \frac{4}{3}=\frac{8}{3}$. This is two thirds of the
      perfectly competitive quantity. The price is then
      $p=10-2\times \frac{8}{3}=\frac{14}{3}=4\frac{2}{3}$.
      Each firm makes profits $\pi=\left( p-AC \right)q=\left(
      \frac{14}{3}-2 \right)\frac{4}{3}=\frac{8}{3}\times
      \frac{4}{3}=\frac{32}{9}=3\frac{5}{9}$.
    }
  \item The market is served by two firms but one firm chooses the quantity to
    produce first. The other firm then observes this quantity and chooses its
    own quantity. The price is then set to clear the market. What quantity does
    each firm produce? What is the price? What profit does each firm make?

    \textcolor{solution-color}{
      First we need to find how the follower will react to the leader's choice
      of quantity. This is exactly how the Cournot player's react to each
      other. Therefore firm 2's reaction function is:
      \[
        q_2\left( q_1 \right)=2-\frac{1}{2}q_1
      \]
      Firm 1 then will take how firm 2 will react into account when deciding on
      a quantity. Using this expression in firm 1's profit function:
      \begin{align*}
        \pi_1\left( q_1 \right) & =  10q_1 - 2q_1^2 -2q_1q_2-2q_1 \\
        \pi_1\left( q_1 \right) &= 8q_1 - 2q_1^2 -2q_1\left(
        2-\frac{1}{2}q_1 \right) \\
        \pi_1\left( q_1 \right) &= 8q_1 - 2q_1^2 -4q_1-q_1^2 \\
        \pi_1\left( q_1 \right) &= 4q_1 - q_1^2  \\
      \end{align*}
      Taking the derivative of this and setting equal to zero:
      \[
        4- 2q_1=0 \;\;\;\;\Longleftrightarrow \;\;\;\; q_1 = 2
      \]
      Note that this is the same as a monopolist quantity. The follower then
      will produce:
      \[
        q_2 = 2-\frac{1}{2}q_1 = 2-\frac{1}{2}\times 2=1
      \]
      This is one quarter of the perfectly competitive quantity. The industry
      as a whole then produces $2+1=3$. This is three quarters of the perfectly
      competitive quantity. The price is then $p=10-2\times3=4$. The profit for
      the leader is $\pi_1=\left( 4-2 \right)\times 2 = 4$. The profit for the
      follower is $\pi_2=\left( 4-2 \right)\times 1 = 2$.
    }

    \textcolor{solution-color}{
      In summary, the answer for each sub-part of this question is:
      \begin{center}
        \begin{tabular}{l|lll} \hline\hline
          & Industry Quantity & Price & Profit for each firm \\ \hline
          Perfect competition &  4 & 2 & 0 \\
          Monopoly & 2 & 6 & 8 \\
          Bertrand & 4 & 2 & 0 \\
          Cournot & $2\frac{2}{3}$ (where $q_1=q_2=1\frac{1}{3}$) & $4\frac{2}{3}$ &
          $3\frac{5}{9}$ \\
          Stackelberg & 3 (where $q_1=2$, $q_2=1$) & 4 & $\pi_1=4$, $\pi_2=2$
          \\ \hline\hline
        \end{tabular}
      \end{center}
    }
\end{subquestions}
\end{document}
